% Part: turing-machines
% Chapter: machines-computations
% Section: turing-machines

\documentclass[../../../include/open-logic-section]{subfiles}

\begin{document}

% SEGMENT OLP-0256-S01

\olfileid{tur}{mac}{tur}
\olsection{டியூரிங் பொறிகள்}

\begin{explain}
ஒரு டியூரிங் பொறியை உருவாக்குவது எது என்பதற்கான முறைசார் வரையறை
கருத்தியல் தன்மையுள்ளதாகத் தோன்றினாலும் உண்மையில் எளியது: ஒரு டியூரிங்
பொறியின் செயல்பாட்டைக் குறிப்பிடத் தேவையான எல்லாத் தகவல்களையும் ஒரே
கணிதக் கட்டமைப்பில் அது பொதிந்து வைக்கிறது. இதில் (1) பொறி இருக்கக்கூடிய
நிலைகள் எவை, (2) நாடாவில் இருக்க அனுமதிக்கப்படும் குறிகள் எவை, (3) பொறி
எந்த நிலையில் தொடங்க வேண்டும், (4) பொறியின் கட்டளைத் தொகுதி என்ன ஆகியவை
அடங்கும்.
\end{explain}

\begin{defn}[டியூரிங் பொறி]
ஒரு \emph{டியூரிங் பொறி} $M$ என்பது பின்வருவனவற்றைக் கொண்ட
$\langle Q, \Sigma, q_0, \delta\rangle$ என்ற நான்குறுப்பு:
\begin{enumerate}
\item \emph{நிலைகளின்} ஒரு முடிவுறு கணம்~$Q$,
\item $\TMendtape$, $\TMblank$ ஆகியவற்றை உள்ளடக்கிய ஒரு முடிவுறு
  \emph{குறியெழுத்துத் தொகுதி} $\Sigma$,
\item ஓர் \emph{தொடக்க நிலை}~$q_0 \in Q$,
\item ஒரு முடிவுறு \emph{கட்டளைத் தொகுதி}~$\delta\colon Q \times \Sigma
  \pto Q \times \Sigma \times \{\TMleft, \TMright, \TMstay\}$.
\end{enumerate}
பகுதிச் சார்பு~$\delta$ ஆனது~$M$ இன் \emph{நிலைமாற்றுச் சார்பு}
என்றும் அழைக்கப்படுகிறது.
\end{defn}

\begin{explain}
நாடா ஒரு திசையில் மட்டுமே முடிவின்றி நீள்கிறது என்று நாம் கருதுகிறோம்.
இதனால் நாடாவின் இடது முனையைக் குறிப்பதற்கான சிறப்புக் குறியாக
~$\TMendtape$ ஐ ஒதுக்குவது பயனுள்ளது. டியூரிங் பொறி நிரல்கள் நாடாவை
விட்டு வெளியேறிவிடும் ``அபாயத்தில்'' எப்போது இருக்கின்றன என்பதை அறிய
இது உதவுகிறது. இக்குறி ஒருபோதும் மேலெழுதப்படாது, அதாவது,
$\delta(q,\TMendtape)$ வரையறுக்கப்பட்டிருந்தால்
$\delta(q,\TMendtape) = \tuple{q', \TMendtape, x}$ என்று நாம்
கருதலாம். சில பாடநூல்கள் இவ்வாறு செய்கின்றன; நாம் அவ்வாறு செய்வதில்லை.
உங்கள் டியூரிங் பொறியைக் கட்டமைக்கும்போது அது~$\TMendtape$ ஐ ஒருபோதும்
மேலெழுதாதபடி வெறுமனே கவனமாக இருக்கலாம். மேலும், அத்தகைய மேலெழுதலை
அனுமதிப்பது வசதியான சில நெகிழ்வை வழங்கும் நேர்வுகளும் உள்ளன.
\end{explain}

\begin{ex}
\emph{இரட்டை எண் பொறி:} இரட்டை எண் பொறி முறைசார்ந்து
$\tuple{Q, \Sigma, q_0, \delta}$ என்ற நான்குறுப்பு; இங்கு
\begin{align*}
Q & = \{ q_0, q_1 \} \\
\Sigma & = \{ \TMendtape, \TMblank, \TMstroke \}, \\
\delta(q_0, \TMstroke) & = \tuple{q_1, \TMstroke, \TMright},\\
\delta(q_1, \TMstroke) & = \tuple{q_0, \TMstroke, \TMright},\\
\delta(q_1, \TMblank)  & = \tuple{q_1, \TMblank, \TMright}.
\end{align*}
\end{ex}

\end{document}
